\documentclass{article}
\usepackage[T1]{fontenc}
\usepackage{lmodern}
\usepackage{graphicx}
\usepackage{amsmath,amssymb}
\usepackage[a4paper,margin=2cm]{geometry}
\usepackage[absolute,overlay]{textpos}
\setlength{\TPHorizModule}{1mm}
\setlength{\TPVertModule}{1mm}
\usepackage[indonesian]{babel}
\usepackage{hyperref}
\setlength{\emergencystretch}{2em}
% unit: Halaman 20

\begin{document}

% === Halaman 20 (llm) ===

Contoh \textit{Null Cipher} lainnya:

\textcolor{blue}{B}ig \textcolor{blue}{r}umble in \textcolor{blue}{N}ew \textcolor{blue}{G}uinea.\\
\textcolor{blue}{T}he \textcolor{blue}{w}ar \textcolor{blue}{o}n \textcolor{blue}{c}elebrity \textcolor{blue}{a}cts \textcolor{blue}{s}hould \textcolor{blue}{e}nd \textcolor{blue}{s}oon.\\
\textcolor{blue}{O}ver \textcolor{blue}{f}our \textcolor{blue}{d}ie \textcolor{blue}{e}cstatic \textcolor{blue}{e}lephants \textcolor{blue}{r}eplicated.

\vspace{5mm}

\textcolor{red}{Bring two cases of deer.}

\end{document}
